% from main.tex
%%%%%%%%%%  scan idx 147  |  folio 134  %%%%%%%%%%
plate, et le fait que \(\overset{\mathrm L}{\otimes}_A\) commute au passage aux fibres).

\textbf{Exercices 5.7.} \textbf{a)} Supposons \(A\) commutatif et \((S',A')\) fidèlement plat sur \((S,A)\) (cf. (2.16.2)). Pour que \(E\in\Ob\Dpx^-({}_A S)\)\footnote{\(S\) = objet final de \(S\)} soit d'amplitude plate contenue dans \([m,n]\), il faut et il suffit que \(f^*(E)\) le soit.

\textbf{b)} Soit \(F\in\Ob\Dpx^+({}_A S)\)\footnotemark[\value{footnote}]. Montrer l'équivalence des conditions :
\begin{enumerate}[label=(\roman*)]
\item \(F\) est isomorphe, dans \(\Dpx^+({}_A S)\), à un complexe \(F'\) tel que \(F'^i\) soit injectif pour tout \(i\) et nul pour \(i\notin[m,n]\);
\item pour tout \(A_S\)-Module à gauche \(E\), on a \(\Ext^i(E,F)=0\) pour \(i\notin[m,n]\).
\end{enumerate}
Si \(F\) satisfait aux conditions précédentes, on dit que \(F\) est d'amplitude injective contenue dans \([m,n]\). La sous-catégorie pleine de \(\Dpx^+({}_A S)\) formée des objets d'amplitude injective finie est une sous-catégorie triangulée, notée \(\Dpx^+({}_A S)_{\injf}\). Montrer que, si \(A\) est commutatif, le foncteur
\[
  \mathrm R\Hom:\Dpx(S)^\circ\times\Dpx^+(S)\longrightarrow \Dpx(S)
\]
induit un foncteur
\[
  \mathrm R\Hom:\Dpx(S)^\circ_{\torf}\times\Dpx(S)_{\injf}
  \longrightarrow \Dpx(S)_{\injf}.
\]

Nous pouvons maintenant répondre à la question de (5.0) et énoncer le résultat principal de ce numéro.

\textbf{Théorème 5.8.} Soient \((S,A)\) un topos annelé, \(E\in\Ob\Dpx({}_A S)\), \(m,n\in\mathbb Z\), avec \(m\leq n\). Conditions équivalentes :
\begin{enumerate}[label=(\roman*)]
\item \(\perfamp(E)\subset[m,n]\) (4.8);
\item \(E\) est \((m-1)\)-pseudo-cohérent (2.15) et l'on a \(\toramp(E)\subset[m,n]\) (5.2).
\end{enumerate}
Comme « parfait » implique évidemment « pseudo-cohérent », on en déduit :

